\SourceLeaf{163}wo
\[
\begin{array}{rl}
A &= 0{,}48213\\
\log B &= 6{,}66231 - 10\\
\log C &= 9{,}67925 - 10\\
\log b &= 0{,}03909.7 \\
\log c &= -0{,}00422.25
\end{array}
\]


Der Werth von $A$ ergibt sich aus
\[
\begin{array}{r@{\quad}l|r@{\quad}l}
n=3 & 0{,}46347 & n=52 & 0{,}48081\\
7 & 0{,}48097 & 57 & 0{,}48409\\
12 & 0{,}48459 & 62 & 0{,}48527\\
17 & 0{,}47976 & 67 & 0{,}48097\\
22 & 0{,}47966 & 72 & 0{,}47736\\
27 & 0{,}48097 & 77 & 0{,}48150\\
32 & 0{,}48224 & 82 & 0{,}49766\\
37 & 0{,}48463 & 87 & 0{,}48097\\
42 & 0{,}48325 & 92 & 0{,}47698\\
47 & 0{,}48097 & 97 & 0{,}51417
\end{array}
\]

\Needspace{7\baselineskip}
\begin{center}
BEMERKUNGEN.
\end{center}

Die Aufgaben [I.] und [II.] sowie die erste Notiz unter [V.], S.\,149/150, sind ein und demselben Handbuch entnommen. [VI.] befindet sich auf einem Blatte am Ende des Buchs der \textsc{Gauss}-Bibliothek: Essai sur les probabilités de la durée de la vie humaine. Par M. \textsc{Deparcieux}. Paris, 1746.

\hfill \textsc{Börsch}, \textsc{Krüger}.
